\chapter{Endre Szemer\'edi (2012)}

\section{Biographical Sketch}

Endre Szemer\'edi was born on 21 August 1940 in Budapest, Hungary. He studied at E\"otv\"os Lor\'and University in Budapest, where he completed his PhD in 1970 under the supervision of Paul Erd\H{o}s. His doctoral work on extremal graph theory and arithmetic progressions immediately established him as one of the leading combinatorialists of his generation.

Szemer\'edi held positions at the Hungarian Academy of Sciences (1965--1973), the University of Chicago (1973--1978), and Rutgers University (1978--present), where he is now professor emeritus. He has received the Abel Prize in 2012, the Wolf Prize in Mathematics in 1992, the Leroy P. Steele Prize in 2008, and the Rolf Schock Prize in 2008.

Szemer\'edi was a frequent collaborator with Paul Erd\H{o}s and was known for his deep, persistent approach to solving fundamental problems. His most celebrated result---Szemer\'edi's theorem on arithmetic progressions---took several years of sustained effort and required the development of an entirely new technique, the Szemer\'edi regularity lemma, which has become one of the most powerful and widely applied tools in combinatorics.

\section{High-Level View for a General Audience}

\subsection{What Did Szemer\'edi Do?}

Imagine you have a large collection of numbers---say, the first million positive integers. Now randomly select 1\% of them (10,000 numbers). Will this random selection necessarily contain three equally spaced numbers (like 17, 22, 27)?

Surprisingly, the answer is no: a random selection of 1\% of the integers typically does not contain a three-term arithmetic progression if you only look at numbers up to some bound. But what if the selection is not random? What if you select a fixed fraction of the numbers?

Endre Szemer\'edi proved a stunning result: any subset of the integers with positive density (that is, any set that contains at least $\delta N$ of the first $N$ integers, for some fixed $\delta > 0$ and all sufficiently large $N$) contains arithmetic progressions of arbitrarily great length.

This is deeply counterintuitive. A random subset of density $\delta$ has expected number $O(\delta^3 N^2)$ of three-term progressions, which goes to zero relative to $N$ when $\delta$ is small. Yet Szemer\'edi's theorem asserts that \emph{any} subset, no matter how structured, must contain progressions once $N$ is large enough relative to $\delta$ and $k$.

The theorem completely settled a 1936 conjecture of Erd\H{o}s and Tur\'an.

\subsection{The Core Ideas}

Szemer\'edi's theorem is proved by an ingenious combinatorial argument that can be summarized as follows:
\begin{enumerate}
\item Suppose $A \subset \{1,\ldots,N\}$ has density $\delta$ but contains no $k$-term arithmetic progression.
\item Show that $A$ must be ``irregularly distributed''---its density varies significantly across different parts of $\{1,\ldots,N\}$.
\item Use this irregularity to find a long subprogression $P \subset \{1,\ldots,N\}$ where $A$ has density at least $\delta + \epsilon$ for some $\epsilon > 0$ that depends only on $\delta$ and $k$.
\item Repeat: we obtain a nested sequence of progressions, each with density increased by $\epsilon$. Since density cannot exceed 1, the process must terminate after finitely many steps.
\item Therefore, for any initial $\delta > 0$ and $k \geq 3$, there exists $N(k,\delta)$ such that any subset of $\{1,\ldots,N\}$ of size $\delta N$ contains a $k$-term progression.
\end{enumerate}

This density increment strategy, while simple in outline, is extremely difficult to implement. The main technical challenge is step 2-3: proving that a progression-free set must have a density increment on some subprogression.

\subsection{Why Does It Matter?}

Szemer\'edi's theorem is one of the deepest results in combinatorics:
\begin{itemize}
\item It resolved a 40-year-old conjecture of Erd\H{o}s and Tur\'an.
\item Its proof introduced the Szemer\'edi regularity lemma, one of the most powerful tools in graph theory.
\item It led to the development of additive combinatorics and higher-order Fourier analysis.
\item It inspired the Green--Tao theorem that primes contain arbitrarily long arithmetic progressions.
\item The regularity lemma has found applications in computer science, including property testing and machine learning.
\end{itemize}

\section{The Origin of the Problem}

\subsection{Van der Waerden's Theorem}

In 1927, Bartel Leendert van der Waerden proved that if the integers are colored with finitely many colors, then one color contains arbitrarily long arithmetic progressions. This was a key result in Ramsey theory.

Van der Waerden's theorem states: for any positive integers $r$ and $k$, there exists $W(r,k)$ such that any $r$-coloring of $\{1,\ldots,W(r,k)\}$ contains a monochromatic $k$-term arithmetic progression.

The numbers $W(r,k)$---the van der Waerden numbers---grow extremely quickly. Even $W(2,3) = 9$ is small, but $W(2,4) = 35$, $W(2,5) = 178$, $W(2,6) = 1132$, and the growth continues super-exponentially. The best known upper bound, due to Gowers, is:
\[
W(r,k) \leq 2^{2^{r^{2^{2^{k+9}}}}},
\]
a tower of exponentials of height roughly $k$.

Van der Waerden's theorem is a special case of Szemer\'edi's theorem. The coloring can be thought of as a partition of the integers, and at least one color class has density at least $1/r$. Szemer\'edi's theorem vastly generalizes this by showing that any set of positive density---not just one color class---contains long progressions.

\subsection{The Erd\H{o}s--Tur\'an Conjecture}

In 1936, Paul Erd\H{o}s and Paul Tur\'an made a much stronger conjecture: any subset of the positive integers with positive upper density contains arbitrarily long arithmetic progressions.

The upper density of a set $A \subset \mathbb{N}$ is defined as:
\[
\bar{d}(A) = \limsup_{N \to \infty} \frac{|A \cap \{1,\ldots,N\}|}{N}.
\]

Erd\H{o}s and Tur\'an were led to this conjecture by their study of the distribution of primes and by van der Waerden's theorem. They observed that van der Waerden's theorem is a corollary of the conjecture, since any finite coloring partitions $\mathbb{N}$ into at least one color class with positive upper density.

The conjecture was known to be true for $k=3$ (three-term progressions) by a theorem of Roth (1953), but progress on longer progressions was slow. Roth's theorem was proved using Fourier analysis: if $A \subset \{1,\ldots,N\}$ has no 3-term progression, then:
\[
|A| \ll \frac{N}{\log\log N}.
\]

Szemer\'edi proved the conjecture for $k=4$ in 1969, then for general $k$ in 1975. His 1975 paper is one of the masterpieces of twentieth-century combinatorics.

\subsection{Progress Before Szemer\'edi}

Before Szemer\'edi's breakthrough, the following partial results were known:
\begin{itemize}
\item \textbf{Roth's theorem (1953):} Any subset of $\{1,\ldots,N\}$ with density at least $c/\log\log N$ contains a 3-term progression. The full density version (any $\delta > 0$) follows from this with $N$ large enough.
\item \textbf{Szemer\'edi (1969):} Any subset of $\{1,\ldots,N\}$ with density at least $\delta$ contains a 4-term progression.
\item \textbf{Roth (1970):} A simplified proof of Szemer\'edi's result for $k=4$ using Fourier analysis.
\end{itemize}

The gap between $k=4$ and general $k$ was enormous. The Fourier-analytic methods used for $k=3$ and $k=4$ break down completely for $k \geq 5$ because the relevant equations are not quadratic.

Szemer\'edi's breakthrough was to develop an entirely new approach based on graph theory and the regularity lemma.

\section{Deep Insight into the Problem}

\subsection{The Connection Between Progressions and Graphs}

Szemer\'edi's brilliant idea was to encode arithmetic progressions in terms of graph properties. For $k=3$ (three-term progressions), the connection is straightforward: a three-term progression $x, x+d, x+2d$ corresponds to a triangle in a 3-uniform hypergraph. For longer progressions, Szemer\'edi used a more sophisticated encoding.

The key insight is that the problem of finding $k$-term arithmetic progressions in a dense set can be reduced to finding certain patterns in graphs and hypergraphs. This reduction allows combinatorial techniques---in particular, the regularity lemma---to be brought to bear.

\subsection{The Density Increment Strategy}

Szemer\'edi's proof uses a density increment strategy:
\begin{enumerate}
\item Assume $A \subset \{1,\ldots,N\}$ has density $\delta$ and contains no $k$-term arithmetic progression.
\item Show that $A$ must be ``irregular'' in some sense: it is not evenly distributed across $\{1,\ldots,N\}$.
\item Using the irregularity, find a long subprogression $P$ of length at least $c_k N^{1/2^k}$ (roughly) where $A$'s density is at least $\delta + \epsilon(\delta, k)$.
\item Iterate: after finitely many steps, the density would exceed 1, a contradiction.
\end{enumerate}

The core difficulty is step 3: proving that a progression-free set has a density increment. Szemer\'edi achieved this by transforming the progression problem into a problem about graphs, applying the regularity lemma to these graphs, and using the structural information given by the regularity lemma to locate the density increment.

\subsection{The Regularity Lemma}

The Szemer\'edi regularity lemma is one of the most important tools in combinatorics. It states:

For any $\epsilon > 0$, there exists $M(\epsilon)$ such that every graph $G$ can be partitioned into $k \leq M$ parts such that all but $\epsilon k^2$ of the pairs of parts are $\epsilon$-regular.

A pair $(V_i, V_j)$ is $\epsilon$-regular if for any subsets $A \subset V_i, B \subset V_j$ with $|A| \geq \epsilon |V_i|, |B| \geq \epsilon |V_j|$, the density of edges between $A$ and $B$ differs from the density between $V_i$ and $V_j$ by at most $\epsilon$.

\subsection{Understanding the Regularity Lemma}

The lemma says that any graph can be approximated by a random-like graph, up to a controlled error. The partition is found by an iterative refinement procedure:
\begin{enumerate}
\item Start with an arbitrary partition.
\item If many pairs are not $\varepsilon$-regular, there exist subsets $A \subset V_i, B \subset V_j$ witnessing the irregularity.
\item Refine the partition by splitting parts based on these subsets.
\item Repeat until the partition is $\varepsilon$-regular.
\end{enumerate}

The bound $M(\epsilon)$ obtained from this procedure is enormous---a tower of twos of height roughly $\epsilon^{-5}$. This tower-type bound is unavoidable: Gowers showed that there exist graphs where any $\epsilon$-regular partition must have at least a tower of twos of height $\Omega(\epsilon^{-1/16})$ parts.

Despite this astronomical bound, the regularity lemma is used as a theoretical tool rather than a practical algorithm. It guarantees the existence of a regular partition, which can then be used to deduce combinatorial properties of the graph.

\subsection{The Regularity Lemma in Action}

Once a graph is partitioned into $\epsilon$-regular parts, we can form the \emph{reduced graph} $R$ whose vertices are the parts $V_1,\ldots,V_k$ and whose edges correspond to $\epsilon$-regular pairs with density above some threshold $\beta$.

The reduced graph inherits many properties of the original graph. A fundamental result is the \emph{embedding lemma}: if the reduced graph contains a certain subgraph $H$, then the original graph contains many copies of $H$ as subgraphs.

This combination---the regularity lemma plus the embedding lemma---is the key tool for proving extremal graph theory results. For Szemer\'edi's theorem, the connection is:
\begin{itemize}
\item Construct a graph whose vertices are the numbers $\{1,\ldots,N\}$ and whose edges encode certain additive relations.
\item Show that if $A$ contains no $k$-term progression, then the graph lacks certain substructures.
\item Apply the regularity lemma to deduce that the graph must be highly structured, leading to a density increment for $A$ on some subprogression.
\end{itemize}

\subsection{Different Proofs of Szemer\'edi's Theorem}

Szemer\'edi's theorem now has several different proofs:
\begin{itemize}
\item \textbf{Combinatorial (Szemer\'edi, 1975):} Uses the density increment strategy and the regularity lemma. The original proof is intricate but elementary, requiring no analysis beyond basic graph theory.
\item \textbf{Ergodic\index{Ergodic theory} (Furstenberg, 1977):} Uses the correspondence principle to convert to a multiple recurrence theorem in ergodic theory\index{Ergodic theory}. Furstenberg's proof showed that Szemer\'edi's theorem is equivalent to a statement about measure-preserving systems:
\[
\liminf_{N \to \infty} \frac{1}{N}\sum_{n=1}^N \mu(A \cap T^{-n}A \cap \cdots \cap T^{-(k-1)n}A) > 0
\]
for any set $A$ of positive measure in a probability space with an invertible measure-preserving transformation $T$.
\item \textbf{Fourier Analytic (Gowers, 2001):} Uses higher-order Fourier analysis and inverse theorems for uniformity norms. Gowers introduced the $U^k$ uniformity norms and proved that if a function has large $U^k$-norm, then it correlates with a polynomial phase function.
\item \textbf{Hypergraph (Gowers, 2005; Rodl--Skokan, 2005):} Uses the hypergraph regularity lemma to prove a hypergraph version of Szemer\'edi's theorem, which implies the original theorem as a special case.
\end{itemize}

Each of these proofs has led to important developments in its respective field.

\section{Biggest Contributions and New Tools}

\subsection{Szemer\'edi's Theorem}

Szemer\'edi's theorem (1975): For any $\delta > 0$ and integer $k \geq 3$, there exists $N(k, \delta)$ such that any subset $A \subset \{1, \ldots, N\}$ with $|A| \geq \delta N$ contains a $k$-term arithmetic progression.

The quantitative bounds obtained from the original proof are enormous:
\[
N(k, \delta) \leq \text{tower}(\delta^{-C_k}),
\]
where $\text{tower}(m)$ is a tower of twos of height $m$. Gowers' proof improved this significantly:
\[
N(k, \delta) \leq 2^{2^{\delta^{-c_k}}}
\]
for some constant $c_k$ depending only on $k$. For $k=3$, the best known bound is:
\[
N(3, \delta) \leq 2^{2^{O(\delta^{-2})}},
\]
due to Bourgain and later improved by Bloom and Sisask.

\subsection{The Regularity Lemma}

The regularity lemma has become one of the most widely used tools in combinatorics:
\begin{itemize}
\item It provides a ``rough'' description of any graph as a union of random-like bipartite graphs.
\item It is essential for the proof of Szemer\'edi's theorem and many other results.
\item It has applications to extremal graph theory, Ramsey theory, and algorithms.
\end{itemize}

The lemma has been extended and refined in many ways:
\begin{itemize}
\item \textbf{The strong regularity lemma:} Allows simultaneous regularization of a graph and its complement.
\item \textbf{The algorithmic regularity lemma:} Provides a polynomial-time algorithm for finding a regular partition, at the cost of slightly weaker regularity.
\item \textbf{Sparse regularity lemmas:} Versions of the lemma for sparse graphs, using different notions of regularity.
\item \textbf{Hypergraph regularity lemmas:} Generalizations to $k$-uniform hypergraphs, which are essential for proving hypergraph versions of Szemer\'edi's theorem.
\end{itemize}

\subsection{The Corners Theorem}

Szemer\'edi proved that any subset of $\mathbb{Z}_N^2$ of positive density contains an axis-aligned right triangle (a ``corner''): points $(x, y), (x+d, y), (x, y+d)$ with $d > 0$.

This theorem is important for several reasons:
\begin{itemize}
\item It is equivalent to Szemer\'edi's theorem for $k=3$ (three-term progressions) via a simple transformation.
\item It has a particularly clean proof using the Fourier transform.
\item It generalizes to higher dimensions, where it becomes the density Hales--Jewett theorem.
\item It is closely related to the coloring Hales--Jewett theorem, which is a combinatorial statement about $n$-dimensional tic-tac-toe.
\end{itemize}

\subsection{Graph Removal Lemma}

Szemer\'edi's graph removal lemma states that any graph on $n$ vertices that contains $o(n^v)$ copies of a fixed subgraph $H$ (with $v$ vertices) can be made $H$-free by removing $o(n^2)$ edges.

More precisely: for any graph $H$ and any $\epsilon > 0$, there exists $\delta > 0$ such that if a graph $G$ on $n$ vertices contains at most $\delta n^v$ copies of $H$, then $G$ can be made $H$-free by removing at most $\epsilon n^2$ edges.

The graph removal lemma has many important consequences:
\begin{itemize}
\item \textbf{Property testing:} It provides the theoretical foundation for testing whether a graph is $H$-free by sampling only a constant number of vertices.
\item \textbf{Extremal graph theory:} It implies that the Tur\'an number $\text{ex}(n, H)$---the maximum number of edges in an $H$-free graph on $n$ vertices---is $o(n^2)$ for any non-bipartite $H$.
\item \textbf{Arithmetic combinatorics:} It implies Roth's theorem on 3-term progressions and Szemer\'edi's theorem for $k=3$.
\end{itemize}

\subsection{The Erd\H{o}s--Szemer\'edi Theorem}

Szemer\'edi and Erd\H{o}s proved a fundamental result about the sum-product phenomenon: for any finite set $A \subset \mathbb{Z}$, either $A+A$ or $A \cdot A$ is large. More precisely:
\[
\max\{|A+A|, |A\cdot A|\} \geq c|A|^{1+\epsilon}
\]
for some absolute constants $c, \epsilon > 0$.

This theorem initiated the study of sum-product phenomena, which has become a major area of additive combinatorics. The best known bound is due to Konyagin and Rudnev:
\[
|A \cdot A + A \cdot A| \gg |A|^{4/3}
\]
for $A \subset \mathbb{R}$.

\subsection{The Erd\H{o}s--Szemer\'edi Theorem on Distinct Distances}

In a separate collaboration, Szemer\'edi and Erd\H{o}s studied the problem of distinct distances in the plane. They proved that any set of $n$ points in the plane determines at least $cn^{1/2}$ distinct distances. This was later dramatically improved by Guth and Katz, who proved that at least $c n / \log n$ distinct distances are determined.

\subsection{Graph Theory and Algorithms}

Szemer\'edi's work on graph algorithms, with Ajtai, Koml\'os, and others, produced several fundamental results:
\begin{itemize}
\item \textbf{The Ajtai--Koml\'os--Szemer\'edi theorem:} A dense graph with minimum degree $\delta n$ contains a Hamiltonian cycle.
\item \textbf{The Ajtai--Koml\'os--Szemer\'edi sorting network:} A construction of optimal-depth sorting networks, a fundamental result in computer science.
\item \textbf{The Ajtai--Koml\'os--Szemer\'edi theorem on the Erd\H{o}s--Szekeres problem:} An upper bound on the Ramsey number $R(3,k)$.
\end{itemize}

\section{Impact of the Discovery}

\subsection{Impact on Mathematics}

\begin{itemize}
\item \textbf{Additive Combinatorics}: Szemer\'edi's theorem founded the field of additive combinatorics. The combination of combinatorial, ergodic\index{Ergodic theory}, and Fourier-analytic methods has produced a rich theory with connections to many areas of mathematics.
\item \textbf{Extremal Graph Theory}: The regularity lemma transformed the field. It is now a standard tool for proving results in extremal graph theory, Ramsey theory, and graph coloring.
\item \textbf{Ergodic Theory}: Furstenberg gave a different proof using ergodic theory, opening new connections. This led to the development of ergodic Ramsey theory and the discovery of deep connections between combinatorics and dynamical systems.
\item \textbf{Number Theory}: The Green--Tao theorem on primes builds on Szemer\'edi's theorem. Green and Tao proved that the primes contain arbitrarily long arithmetic progressions by combining Szemer\'edi's theorem with a transference principle that allows passage from dense sets to sparse pseudorandom sets.
\item \textbf{Fourier Analysis}: Gowers' proof using higher-order Fourier analysis opened new directions. The inverse theorems for uniformity norms are fundamental tools in additive combinatorics.
\item \textbf{Hypergraph Theory}: The hypergraph regularity lemma, developed by Gowers, Rodl, Skokan, and others, extends the regularity method to higher-uniformity hypergraphs and has led to new results in extremal hypergraph theory.
\end{itemize}

\subsection{Impact on Computer Science}

\begin{itemize}
\item \textbf{Property Testing}: The regularity lemma is used for testing graph properties. A property tester is an algorithm that examines a small random sample of the input and decides whether the input has a given property or is far from having it. The regularity lemma implies that many graph properties are testable in constant time.
\item \textbf{Data Mining}: Regular partitions are used for clustering and community detection. The notion of a regular partition provides a rigorous foundation for clustering algorithms based on graph structure.
\item \textbf{Algorithm Design}: The algorithmic version of the regularity lemma is used in graph algorithms. While the original regularity lemma gives astronomical bounds, practical algorithms for computing approximate regular partitions have been developed.
\item \textbf{Sorting Networks}: The Ajtai--Koml\'os--Szemer\'edi sorting network is a construction of comparator networks that sort $n$ elements in $O(\log n)$ depth, which is optimal up to constant factors.
\end{itemize}

\section{Current Developments}

\subsection{Quantitative Bounds}

A major direction of current research is improving the bounds. Gowers' proof gave a bound of the form $N \leq 2^{2^{1/\delta^{c_k}}}$, which is much better than the tower-type bounds from the original proof.

Recent work by Peluse, Sah, and others has improved these bounds for $k=3$ and $k=4$:
\begin{itemize}
\item For $k=3$ (Roth's theorem), the best current bound is:
\[
N(3,\delta) \leq 2^{2^{O(\delta^{-1})}},
\]
due to Bloom and Sisask, building on earlier work of Bourgain.
\item For $k=4$, the best bound is due to Green and Tao, who proved that any subset of $\{1,\ldots,N\}$ of size at least $c N (\log N)^{-c}$ contains a 4-term progression. More recently, Peluse proved a polynomial bound for $k=4$ in finite fields.
\end{itemize}

For general $k$, the best bound remains Gowers' bound, and improving it is a major open problem.

\subsection{Higher-Order Fourier Analysis}

Green, Tao, and Ziegler developed the theory of nilsequences and inverse theorems for the Gowers uniformity norms. The central result is the inverse theorem for the $U^k$-norm: if $f: \mathbb{Z}_N \to \mathbb{C}$ is a bounded function with $\|f\|_{U^k} \geq \delta$, then there exists a nilsequence $\psi: \mathbb{Z}_N \to \mathbb{C}$ of complexity bounded by $C(k,\delta)$ such that:
\[
|\langle f, \psi \rangle| \geq c(k,\delta).
\]

This theorem is the key ingredient in the Green--Tao--Ziegler proof of the inverse theorem and has been used to prove many results in additive combinatorics, including:
\begin{itemize}
\item The Mobius and nilsequences conjecture, which relates the Mobius function to nilsequences.
\item The Chowla conjecture on correlations of the Mobius function.
\item The logarithmically averaged versions of the Chowla and Elliott conjectures.
\end{itemize}

\subsection{The Transference Principle}

The transference principle, developed by Green and Tao, extends Szemer\'edi-type results to sparse pseudorandom sets. The idea is:
\begin{enumerate}
\item Let $\nu: \mathbb{Z}_N \to \mathbb{R}_{\geq 0}$ be a ``pseudorandom measure'' that models a sparse set (like the primes, where $\nu(n) \approx \log n$ for primes and 0 otherwise).
\item Suppose $f: \mathbb{Z}_N \to \mathbb{R}_{\geq 0}$ satisfies $0 \leq f \leq \nu$ and $\sum f \geq \delta \sum \nu$.
\item Show that $f$ can be ``transferred'' to a function $\tilde{f}$ on $\mathbb{Z}_N$ with $0 \leq \tilde{f} \leq 1$ and $\int \tilde{f} \geq \delta/2$ such that the number of $k$-term progressions in $f$ is approximately the same as in $\tilde{f}$.
\item Apply Szemer\'edi's theorem to $\tilde{f}$ to conclude that $f$ contains a $k$-term progression.
\end{enumerate}

This was the key idea in the proof that the primes contain arbitrarily long arithmetic progressions. It has since been applied to many other problems in additive combinatorics.

\subsection{Algorithmic Regularity}

Efficient algorithms for computing Szemer\'edi partitions have been developed by Alon, Fischer, Krivelevich, Szegedy, and others. The algorithmic regularity lemma finds a partition with the following properties in polynomial time:
\begin{itemize}
\item The number of parts is $M(\epsilon)$.
\item All but $\epsilon k^2$ pairs of parts are $\epsilon$-regular.
\end{itemize}

The algorithm works by:
\begin{enumerate}
\item Starting with the trivial partition (one part).
\item If a pair $(V_i, V_j)$ is not $\epsilon$-regular, find subsets $A \subset V_i$, $B \subset V_j$ witnessing the irregularity.
\item Use these subsets to refine the partition.
\item Repeat until the partition is $\epsilon$-regular.
\end{enumerate}

While the worst-case running time is polynomial in $n$, the constant $M(\epsilon)$ grows so fast that the algorithm is impractical for large graphs. However, for many practical applications, much weaker regularity suffices.

\subsection{Graph Limits and Regularity}

The theory of graph limits (graphons) builds on the regularity lemma by providing a limit object for convergent sequences of dense graphs. A \emph{graphon} is a symmetric measurable function $W: [0,1]^2 \to [0,1]$. Every sequence of graphs contains a convergent subsequence whose limit is a graphon.

The connection to the regularity lemma is:
\begin{itemize}
\item The regularity lemma approximates a graphon by a step function with finitely many steps.
\item The compactness of the space of graphons (under the cut metric) is equivalent to the analytic form of the regularity lemma.
\item The theory of graph limits provides a continuous framework for studying large graphs and their properties.
\end{itemize}

\subsection{Hypergraph Regularity and Extremal Problems}

The hypergraph regularity lemma has been applied to extremal hypergraph problems, including:
\begin{itemize}
\item The proof of the Erd\H{o}s--Ginzburg--Ziv theorem.
\item The study of hypergraph Tur\'an numbers.
\item The proof of the hypergraph removal lemma, which implies Szemer\'edi's theorem for progressions and the multidimensional Szemer\'edi theorem.
\end{itemize}

The hypergraph regularity lemma is significantly more complex than the graph regularity lemma. A $k$-uniform hypergraph is partitioned into $k$-tuples of parts, and regularity is defined in terms of the density of $k$-uniform edges across $k$-tuples of parts. The number of parts required is enormous---a tower of twos of height depending on $k$ and $\epsilon$.

\subsection{Additive Combinatorics Beyond Progressions}

The methods developed for Szemer\'edi's theorem have been extended to many other combinatorial patterns:
\begin{itemize}
\item \textbf{Polynomial progressions:} Bergelson and Leibman proved that any set of positive density contains patterns of the form $x, x+p_1(n), \ldots, x+p_k(n)$ for any polynomials $p_1,\ldots,p_k$ with zero constant term.
\item \textbf{Multidimensional progressions:} The multidimensional Szemer\'edi theorem, proved by Furstenberg and Katznelson, states that any subset of $\mathbb{Z}^d$ of positive density contains axis-parallel $k$-dimensional cubes.
\item \textbf{Patterns in finite fields:} The finite field analog of Szemer\'edi's theorem, developed by Bergelson, Tao, and Ziegler, has connections to the theory of mixing and the Gowers norms.
\end{itemize}

\subsection{The Erd\H{o}s--Tur\'an Conjecture on Additive Bases}

The original Erd\H{o}s--Tur\'an conjecture was actually about additive bases---sets $A \subset \mathbb{N}$ such that every sufficiently large integer can be written as a sum of two elements of $A$. The conjecture on progressions was a separate problem that turned out to be more tractable.

The additive bases problem---known as the Erd\H{o}s--Tur\'an conjecture on additive bases---remains open. It asks: if $A \subset \mathbb{N}$ is an asymptotic basis of order 2 (every sufficiently large integer is the sum of two elements of $A$), does $A$ necessarily contain arbitrarily long arithmetic progressions?

\section{Conclusion}

Endre Szemer\'edi's theorem on arithmetic progressions is one of the landmarks of late twentieth-century mathematics. The regularity lemma he introduced has transformed combinatorics and found applications throughout mathematics and computer science.

The theorem brings together combinatorics, ergodic theory, Fourier analysis, and number theory in a remarkable synthesis. Each of the different proofs of Szemer\'edi's theorem has opened new connections between these fields, and the methods developed for the theorem have been extended to a wide range of problems.

Szemer\'edi's work exemplifies the power of combinatorial reasoning at its deepest level. His persistence in attacking a problem that resisted solution for forty years, and his creation of an entirely new method to solve it, place him among the greatest combinatorialists of all time.

The regularity lemma, originally developed as a tool for one specific problem, has become an indispensable part of the mathematician's toolkit. Its applications in property testing, graph limits, and algorithmic graph theory continue to grow.


\section{Behind the Breakthrough: The Human Story}
Szemer\'edi was advised by Tur\'an, who gave him a problem on arithmetic progressions. Szemer\'edi struggled for years before solving it, leading to a theorem that Paul Erd\H{o}s called a masterpiece. Erd\H{o}s had offered a \$1,000 prize for the proof, which he gladly paid.

\section{Pedagogical Metaphors and Lineage}
\subsection{Signature Metaphor}
``Finding order in chaos by proving that any large graph or network can be broken down into a few highly organized, random-like pieces.''

\subsection{Mathematical Lineage}
Advised by Paul Tur\'an, who was advised by Lip\'ot Fej\'er.

\section{Applied Science and Computational Algorithms}
\subsection{Key Takeaways for Applied Science}
Used to analyze large-scale network topologies, such as community detection in social graphs, finding sub-structures in web indexes, and in property testing algorithms in computer science.

\subsection{Algorithms and Numerical Implementations}
Algorithmic versions of Szemer\'edi's Regularity Lemma (such as the Alon-Duke-Lefmann-Rodl-Yuster algorithm) construct $\varepsilon$-regular partitions in polynomial time for graph approximation.

\section{The Research Frontier}
Finding optimal bounds for Szemer\'edi's theorem and developing the hypergraph regularity lemma to solve multi-dimensional patterns in combinatorics.

\section{Guided Exercises and Challenges}
\begin{enumerate}[label=\textbf{\arabic*.}]
  \item State Szemer\'edi's theorem on arithmetic progressions and explain the difference between it and van der Waerden's theorem.
  \item Formulate Szemer\'edi's Regularity Lemma for graphs and define $\varepsilon$-regular pairs.
  \item Describe the Roth's theorem case ($k=3$) and prove it using the triangle removal lemma.
\end{enumerate}

\section{Curriculum Map and Further Reading}
For students wishing to delve deeper into the mathematical machinery underlying these breakthroughs, the following learning path is recommended:
\begin{itemize}
  \item David Conlon, Jacob Fox, and Yufei Zhao, \emph{A Guide to Szemer\'edi's Regularity Lemma}.
  \item Terence Tao and Van H. Vu, \emph{Additive Combinatorics}, Cambridge University Press.
\end{itemize}

\section*{Bibliography}
\begin{enumerate}
\item E. Szemer\'edi, ``On sets of integers containing no $k$ elements in arithmetic progression,'' Acta Arithmetica 27 (1975), 199--245.
\item E. Szemer\'edi, ``Regular partitions of graphs,'' in ``Probl\`emes Combinatoires et Th\'eorie des Graphes,'' CNRS, 1978, 399--401.
\item B. Green and T. Tao, ``The primes contain arbitrarily long arithmetic progressions,'' Annals of Mathematics 167 (2008), 481--547.
\item W. T. Gowers, ``A new proof of Szemer\'edi's theorem,'' Geometric and Functional Analysis 11 (2001), 465--588.
\item H. Furstenberg, ``Ergodic behavior of diagonal measures and a theorem of Szemer\'edi on arithmetic progressions,'' Journal d'Analyse Math\'ematique 31 (1977), 204--256.
\item T. Tao and V. Vu, ``Additive Combinatorics,'' Cambridge University Press, 2006.
\item D. Conlon and J. Fox, ``Graph removal lemmas,'' in ``Surveys in Combinatorics 2013,'' Cambridge University Press, 2013.
\item L. Lov\'asz and B. Szegedy, ``Limits of dense graph sequences,'' Journal of Combinatorial Theory B 96 (2006), 933--957.
\item J. Koml\'os and M. Simonovits, ``Szemer\'edi's regularity lemma and its applications in graph theory,'' in ``Combinatorics, Paul Erd\H{o}s is Eighty,'' Bolyai Society, 1993.
\item J. Koml\'os, A. Shokoufandeh, M. Simonovits, and E. Szemer\'edi, ``The regularity lemma and its applications in graph theory,'' in ``Theoretical Aspects of Computer Science,'' Springer, 2002.
\item R. Graham, B. Rothschild, and J. Spencer, ``Ramsey Theory,'' Wiley, 1990.
\item P. Erd\H{o}s and P. Tur\'an, ``On some sequences of integers,'' Journal of the London Mathematical Society 11 (1936), 261--264.
\item B. L. van der Waerden, ``Beweis einer Baudetschen Vermutung,'' Nieuw Archief voor Wiskunde 15 (1927), 212--216.
\item T. Tao, ``The dichotomy between structure and randomness, arithmetic progressions, and the primes,'' in ``Proceedings of the ICM 2006.''
\item J. Fox and L. M. Lov\'asz, ``A tight bound for Szemer\'edi's regularity lemma,'' Combinatorica 37 (2017), 911--951.
\item K. F. Roth, ``On certain sets of integers,'' Journal of the London Mathematical Society 28 (1953), 104--109.
\item H. Furstenberg and Y. Katznelson, ``An ergodic Szemer\'edi theorem for commuting transformations,'' Journal d'Analyse Math\'ematique 38 (1978), 275--291.
\item B. Green and T. Tao, ``An inverse theorem for the Gowers $U^3$-norm,'' Proceedings of the Edinburgh Mathematical Society 51 (2008), 73--105.
\item B. Green, T. Tao, and T. Ziegler, ``An inverse theorem for the Gowers $U^s$-norm,'' Annals of Mathematics 176 (2012), 1231--1372.
\item T. Gowers, ``Quasirandomness, counting and regularity for 3-uniform hypergraphs,'' Combinatorics, Probability and Computing 15 (2006), 143--184.
\item V. Rodl and J. Skokan, ``Regularity lemma for $k$-uniform hypergraphs,'' Random Structures and Algorithms 25 (2004), 1--42.
\item T. Bloom and O. Sisask, ``Breaking the logarithmic barrier in Roth's theorem on arithmetic progressions,'' arXiv:2007.03528 (2020).
\item S. Peluse, ``Bounds for sets with no polynomial progressions,'' arXiv:1908.07327 (2019).
\item L. Guth and N. Katz, ``On the Erd\H{o}s distinct distances problem in the plane,'' Annals of Mathematics 181 (2015), 155--190.
\item M. Ajtai, J. Koml\'os, and E. Szemer\'edi, ``Sorting in $c \log n$ parallel steps,'' Combinatorica 3 (1983), 1--19.
\item E. Szemer\'edi, ``The origin of the regularity lemma,'' in ``The Abel Prize 2012,'' Springer, 2014.
\item J. Bourgain, ``On triples in arithmetic progression,'' Geometric and Functional Analysis 9 (1999), 968--984.
\item V. Bergelson and A. Leibman, ``Polynomial extensions of van der Waerden and Szemer\'edi theorems,'' Journal of the AMS 9 (1996), 725--753.
\end{enumerate}
